% BHU PYQ -- Manually transcribed & error-corrected
\documentclass[11pt,a4paper]{article}

\usepackage[a4paper,top=2.5cm,bottom=2.5cm,left=2.8cm,right=2.8cm,headheight=14pt]{geometry}
\usepackage{amsmath,amssymb}
\usepackage[T1]{fontenc}
\usepackage[utf8]{inputenc}
\usepackage{lmodern}
\usepackage{microtype}
\usepackage{fancyhdr}
\usepackage{enumitem}
\usepackage{hyperref}
\usepackage{lastpage}
\usepackage{parskip}

\hypersetup{colorlinks=true,urlcolor=black,linkcolor=black,
  pdftitle={CHB-401: Inorganic Chemistry-II (Old Course) -- BHU PYQ},pdfauthor={Banaras Hindu University}}

\pagestyle{fancy}\fancyhf{}
\renewcommand{\headrulewidth}{0.4pt}\renewcommand{\footrulewidth}{0.4pt}
\fancyhead[L]{\small   \textit{Banaras Hindu University}}
\fancyhead[R]{\small   \textit{CHB-401: Inorganic Chemistry-II (Old Course)}}
\fancyfoot[C]{\small Page  \thepage\ of \pageref{LastPage}}


\newlist{parts}{enumerate}{2}
\setlist[parts,1]{label=(\alph*),leftmargin=*,itemsep=5pt,topsep=3pt}
\setlist[parts,2]{label=(\roman*),leftmargin=*,itemsep=3pt,topsep=2pt}

\newcommand{\pts}[1]{\hfill   \textit{[#1]}}


\begin{document}


\begin{center}
  {\large\bfseries BANARAS HINDU UNIVERSITY}\\[2pt]
  {
\normalsize B.Sc. (Hons.) Semester IV Examination, 2013-14 (Old Course)}\\[6pt]
  \rule{0.8\linewidth}{0.5pt}\\[6pt]
  {\large\bfseries Paper No. CHB-401: Inorganic Chemistry-II and Selected Topics}\\[4pt]
  {
\normalsize Time: 3 Hours\hfill Maximum Marks: 70}
\end{center}
\rule{\linewidth}{0.4pt}
\small



\noindent   \textit{Instructions: Answer five questions in total, including Question~1 from each section which is compulsory. Use separate answer books for Section~A and Section~B.}

\normalsize
\bigskip


\begin{center}
  {\large\bfseries SECTION A}\\[2pt]
  {
\normalsize (Inorganic Chemistry-II -- Marks: 35)}
\end{center}
\rule{0.4\linewidth}{0.2pt}
\smallskip




\noindent   \textbf{Question 1.} Answer all parts of the following: \pts{5 $ \times$ 2 = 10}

\begin{parts}
  \item How are pure superoxo and peroxo salts prepared in liquid ammonia?
  \item Why are $ \text{Cu}^{2+}$ complexes colored, whereas $ \text{Cu}^+$ complexes are colorless?
  \item Why are ammine complexes of lanthanide ions not formed in aqueous solutions?
  \item Why is $ \text{NH}_3$ a stronger Lewis base than $ \text{NF}_3$?
  \item Define double salts and coordination compounds with examples.
\end{parts}

\medskip\rule{\linewidth}{0.2pt}




\noindent   \textbf{Question 2.}

\begin{parts}
  \item State the Br\o nsted-Lowry definition of acids and bases. Illustrate Lewis acids and bases with equations. \pts{6}
  \item In terms of HSAB principles, which end of the $ \text{SCN}^-$ (thiocyanate) ion would you expect to coordinate to $ \text{Co}^{2+}$ and $ \text{Cr}^{3+}$? Explain. \pts{6}
\end{parts}

\medskip\rule{\linewidth}{0.2pt}




\noindent   \textbf{Question 3.}

\begin{parts}
  \item Write the IUPAC names of the following coordination compounds: \pts{6}
  
\begin{parts}
    \item $[ \text{Co(NO}_2 \text{)(NH}_3 \text{)}_5]^{2+}$
    \item $[ \text{Cr(en)}_3] \text{Cl}_3$
    \item $ \text{K}_2[ \text{Zn(NCS)}_4]$
  \end{parts}
  \item Depict the structures and write the names of the isomers of $[ \text{Co(en)}_2 \text{Cl}_2]^+$. \pts{6}
\end{parts}

\medskip\rule{\linewidth}{0.2pt}




\noindent   \textbf{Question 4.}

\begin{parts}
  \item Discuss briefly three classes of aprotic solvents with examples. \pts{6}
  \item Write outer electronic configurations of $ \text{Sm}^{2+}$ and $ \text{Tb}^{4+}$ and indicate the number of unpaired electrons in each case. \pts{6}
  \item Why do lanthanide ions show a weaker tendency towards complex formation as compared to transition metal ions? What is characteristic about their coordination numbers? Explain. \pts{6}
\end{parts}


\newpage

\begin{center}
  {\large\bfseries SECTION B}\\[2pt]
  {
\normalsize (Selected Topics in Chemistry -- Marks: 35)}
\end{center}
\rule{0.4\linewidth}{0.2pt}
\smallskip




\noindent   \textbf{Question 5.} Answer all parts of the following: \pts{5 $ \times$ 2 = 10}

\begin{parts}
  \item What is the basic difference between a fuel cell and a battery?
  \item What is biomass? How can it be used as a renewable source of energy?
  \item Discuss the cathodic protection method for the prevention of metallic corrosion.
  \item Give a brief account of photoisomerization involving rotation about a $ \text{C=C}$ double bond.
  \item Mention some of the limitations of the applicability of thermodynamics in biological systems.
\end{parts}

\medskip\rule{\linewidth}{0.2pt}




\noindent   \textbf{Question 6.}

\begin{parts}
  \item Describe the construction and working of the $ \text{H}_2 \text{-O}_2$ fuel cell. Mention its components, cell reactions involved, and maximum theoretical e.m.f. obtainable. \pts{8}
  \item Discuss various methods used for the disposal and recycling of polymer wastes. \pts{6}
\end{parts}

\medskip\rule{\linewidth}{0.2pt}




\noindent   \textbf{Question 7.}

\begin{parts}
  \item Name different types of metallic corrosion, illustrating them in brief. \pts{6}
  \item Describe the weight-loss method for the determination of corrosion rates. \pts{6}
  \item What is rust? Write down the mechanism of rust formation on iron. \pts{2}
\end{parts}

\vfill
\rule{\linewidth}{0.4pt}

\begin{center}\small
  \href{https://bhu-exam.vercel.app/}{Click here to visit our website}\\[4pt]
  \href{https://me-aryan.vercel.app/}{Click here to visit our contributor}\\[4pt]
  \href{https://quashan-me.vercel.app/}{Click here to visit our contributor}
\end{center}

\end{document}
